% =====================================================================================
% Revised Real-Robot Setup + Results + Conclusion for swarm_intelligence.tex
% Changed or new text is RED (\rev{...} as defined for paper_methodology_revised.tex).
%
% Replace everything from the fitness-evolution figure (\begin{figure}[t] ... fig:fitness),
% which directly follows the fig:wind figure, down to the end of the Conclusion section
% (just before \begin{credits}) with the text below.
%
% Upload these figures (from hardware_transfer_test/final/overleaf_summary/figures/) to
% Images/results/ in Overleaf:
%   fitness_curves_plain_seed123.pdf, paper_dist_vs_battery_n20.pdf,
%   paper_battery_awareness_boxplot.pdf, paper_trajectories_stage1_vs_stage2.pdf,
%   n_agents_sweep_distance_battery_plain_seed123.pdf
%   paper_leadership_n20.pdf
% Numbers come from paper_results.py (100 runs), n_agents_sweep_analysis.py (30 runs)
% and paper_leadership.py (30 runs, from leadership_metrics.py's seed-sweep tables).
% =====================================================================================

\subsection{\texorpdfstring{\rev{Real-Robot Setup}}{Real-Robot Setup}}
\label{sec:hardware_setup}
\begingroup\color{red}
For the real-robot experiments we use Thymio-II robots, each carrying a Raspberry Pi that runs
the controller. The controller runs at 2\,Hz, the same rate as the simulation step, because
the Hebbian update is applied once per step and a faster rate would change its learning
dynamics. The robots have no range-and-bearing sensor with the reach of the simulated one:
the Thymio's infrared sensors see only about 12\,cm. We therefore track every robot with an
OptiTrack motion-capture system and compute the same ten inputs as in simulation from the
tracked positions and headings. There is also no wind tunnel and no battery telemetry.
Each robot's battery is therefore virtual: it is computed online with the wind, drag and
drainage model of Section~\ref{sec:EXp_Setup}, from the robots' real tracked positions and
the robot's measured speed. This preserves what the controller responds to: its own battery
level and the shelter it gets from the others. The commanded $(v, \omega)$ is converted into
wheel speeds through a calibrated conversion.

We deploy the controller evolved with the safety mechanisms (Section~\ref{sec:safety}), and
the speed clamp of Eq.~\ref{eq:clamp} is enforced on the robots using the tracked gaps
between them. The de-overlap push has no physical counterpart and is not used. In the first
trials the training thresholds ($g_{\text{out}} = 0.30$\,m, $g_{\text{in}} = 0.05$\,m) led to
deadlocks: tracking noise of several centimetres made neighbouring robots appear to overlap,
which set both of their forward speeds to zero. We therefore loosened the thresholds on the
robots to $g_{\text{out}} = 0.12$\,m and $g_{\text{in}} = -0.03$\,m. The deployed controller
thus brakes differently from how it was trained.
\endgroup
\todo[inline]{Number of robots, arena size, number of trials and trial duration.}


\begin{figure}[t]
    \centering
    \includegraphics[width=\linewidth]{Images/results/fitness_curves_plain_seed123}
    \caption[]{\rev{Fitness evolution during the two evolutionary stages: fitness of the best candidate in each generation, for the controller evolved without (orange) and with (blue) the safety mechanisms. Each stage has its own fitness function, so the two panels are not comparable. Stars mark the best candidate ever evaluated, which is the result of the stage.}}
    \label{fig:fitness}
\end{figure}


\begin{table}[t]
  \centering
  \caption{Simulation parameters for the baseline controller.}
  \label{tab:baseline_params}
  \begin{tabular}{@{}lll@{}}
    \toprule
    \textbf{Parameter} & \textbf{Description} & \textbf{Value} \\
    \midrule
    $K_{1}$     & MDMC linear gain                           & $0.05~\text{m/s}$ \\
    $K_{2}$     & MDMC angular gain                          & $0.5~\text{rad/s}$ \\
    $\alpha$    & Steepness of potential function            & $2$ \\
    $\epsilon$  & Strength of potential function             & $1$ \\
    $d_{\text{des}}$ & Desired inter–robot distance         & $0.7~\text{m}$ \\
    $D_{\text{p}}$       & Max. interaction range (proximal control)  & $3~\text{m}$ \\
    $K_{\text{goal}}$  & Goal seeking weight                        & $3$ \\
    \rev{$B_0$} & \rev{Starting battery} & \rev{$150$} \\
    \bottomrule
  \end{tabular}
\end{table}


\begin{figure}[t]
  \centering
  \subfloat[\label{fig:scatter_batt_dist}]{
    \begin{minipage}[t]{0.55\textwidth}
      \centering
      \includegraphics[width=0.98\linewidth]{Images/results/paper_dist_vs_battery_n20}
    \end{minipage}
  }
  \subfloat[\label{fig:boxplot}]{
    \begin{minipage}[t]{0.40\textwidth}
     \centering
      \includegraphics[width=0.95\linewidth]{Images/results/paper_battery_awareness_boxplot}
    \end{minipage}
  }
\caption{(a) Travelled distance against remaining battery of 100 simulations \rev{of 20 robots with the best controller of stage 1 (1), stage 2 (2), stage 2 evolved with the safety mechanisms (3), and the baseline (4); ellipses show one standard deviation.} Qualitative comparisons of all stages are provided in the supplementary video \cite{mahdavi_nasab_2026_ants}. (b) Wind speed experienced by the agents initialized at 50\% (left), and 100\% (right) over 100 simulations.}
  \label{fig:joined_scatter_box}
\end{figure}


\begin{table}[t]
\centering
\caption{\rev{Performance over 100 simulations of 20 robots (mean $\pm$ standard deviation). ``Both better'' counts the runs in which a controller travels further \emph{and} keeps more battery than the baseline on the same starting configuration.}}
\label{tab:results_main}
{\color{red}
\begin{tabular}{lccc}
\toprule
\textbf{Controller} & \textbf{Distance [m]} & \textbf{Battery [\%]} & \textbf{Both better} \\
\midrule
Stage 1 & $12.32 \pm 0.74$ & $17.87 \pm 4.87$ & 42/100 \\
Stage 2 & $19.78 \pm 2.05$ & $21.25 \pm 4.85$ & 69/100 \\
Stage 2, safety mechanisms & $13.90 \pm 0.67$ & $16.11 \pm 3.49$ & 37/100 \\
Baseline & $8.73 \pm 0.82$ & $18.31 \pm 4.70$ & -- \\
\bottomrule
\end{tabular}}
\end{table}


\begin{table}[t]
\centering
\caption{Results of battery awareness experiment (\rev{$p=5.0\times10^{-10}$}, $df=198$)}
\begin{tabular}{lccc}
\toprule
\textbf{Experiment} & \textbf{No. of runs} & \textbf{Mean} & \textbf{Standard deviation} \\
\midrule
Agent with 100\% battery & 100 & \rev{60.0} & \rev{13.0} \\
Agent with 50\% battery & 100 & \rev{49.0} & \rev{10.8} \\
\bottomrule
\end{tabular}
\label{tab:results_batt_aware}
\end{table}


\section{Experimental Results}
\label{sec:Exp:Results}

The fitness progression of the evolution of the ABCD-rules on a swarm of 20 agents across 300 generations (\rev{100 in stage 1 and 200 in stage 2}, as described in Section~\ref{sec:stages}) is depicted in Fig.~\ref{fig:fitness}. \rev{In stage 2 the best fitness peaks after 115 generations without the safety mechanisms and after 82 generations with them, and declines afterwards. Because we keep the best candidate ever evaluated, the resulting controllers come from these peaks.} After evolution, we evaluate the best controller obtained after each evolutionary stage and compare it with the standard collective motion baseline~\cite{article}. The baseline is implemented using only proximal control and goal-seeking terms. The parameters are reported in Table~\ref{tab:baseline_params}. \rev{The baseline runs in the same simulator, with the same wind, drag and battery model and the same starting configurations. Its robots start with 150 battery units, the native scale of the baseline model, rather than 100. Since the drain per step is the same, they start with $1.5\times$ as much energy; all battery values are reported as a percentage of each controller's own starting charge.}

Each controller is evaluated over 100 simulation runs of a swarm of 20 robots, initialized with random positions and random network weights; each simulation terminates when the battery of any robot is depleted. Performance is measured as the distance travelled to the left at termination and the average remaining battery across the swarm. \rev{None of the evaluation runs uses a starting configuration seen during evolution.}

Fig.~\ref{fig:scatter_batt_dist} \rev{and Table~\ref{tab:results_main}} summarize the main results. Each point corresponds to one simulation of the best controller from stage 1 (cluster 1), stage 2 (cluster 2), \rev{stage 2 evolved with the safety mechanisms (cluster 3)}, or the baseline (cluster 4), plotted in terms of distance travelled versus the average battery level of the swarm members at the end of the simulation. \rev{Because stage 1 already trains with the wind on and rewards progress against it, the stage-1 controller already travels further than the baseline ($12.3$ vs.\ $8.7$\,m) while keeping a similar amount of battery. Stage 2, which additionally rewards remaining battery and penalises collisions, travels $19.8$\,m, $127\%$ further than the baseline, while also keeping more battery ($21.3\%$ vs.\ $18.3\%$). It beats the baseline on both measures at once in 69 of the 100 runs. The controller evolved with the safety mechanisms trades part of this gain for collision safety: it travels $13.9$\,m ($59\%$ further than the baseline) but keeps slightly less battery ($16.1\%$), and beats the baseline on both measures in 37 of the 100 runs.}

\rev{\textbf{Swarm size.} The advantage over the baseline depends strongly on the number of robots (Fig.~\ref{fig:swarm_size}). The baseline covers about 25\,m with up to 8 robots, but its distance falls steadily as the swarm grows, to 8.8\,m at 20 robots. The stage-2 controller behaves the other way round: a single robot covers only about 1\,m, and the distance grows with the swarm size to about 20\,m. The controller overtakes the baseline in distance at 12 robots (at 14 with the safety mechanisms). With 12 or fewer robots it never beats the baseline on distance and battery at the same time in any of 30 runs per swarm size. With 13 robots it does so in 2 of 30 runs, and with 20 robots in 23 of 30. The controllers were evolved with 20 robots only, and they do not transfer to small swarms.}

We argue that the strategies learned in \rev{stage 2} rely primarily on local inter-robot interactions rather than the compass information. We contrast the controllers of stages 1 and \rev{2} by testing them with 1 and 5 agents \rev{(Fig.~\ref{fig:Debunk_Compass})}. \rev{The stage-1 controller heads upwind whether it is alone or in a group: a single robot travels about 14\,m, heading straight for the wall and then following it upwind (Fig.~\ref{fig:Debunk_Compass}, top left), and five robots move upwind spread out, each on its own path. The stage-2 controller behaves very differently alone and in a group. A single robot turns in continuous circles and drifts only about a metre (top right), whereas five robots stay together, move upwind and repeatedly loop around one another, changing positions within the group (bottom right). The stage-2 controller therefore depends on its neighbours to make progress at all: its behaviour is collective, arising from local interactions rather than from the compass alone.}

So far, we have shown that the evolved controller achieves energy-efficient flocking with an emergent formation reconfiguration strategy and improved distance-battery trade-offs. However, these results alone do not reveal whether the controller actually exploits the internal battery state. To test whether battery awareness has any impact on the strategy that robots adopt, we compare the behaviour of an agent initialised normally with $100\%$ battery to an agent initialised with $50\%$ battery at the beginning of the simulation. The percentage wind speed of the grid cells where the agent travels, $U_{i,j}^\%$, is averaged over the simulation time. This parameter is chosen to quantify conservativeness; an agent that drives a lot in the front of the swarm will expose itself to much higher average $U_{i,j}^\%$ than an agent that tries to stay in the middle of the swarm shielded by others. Once again, we note that the robots do not have any onboard sensors for wind measurement. We ran 100 simulations with all agents initialized at full charge (100\%) and recorded the average $U_{i,j}^\%$ along the trajectory of a representative agent. We then repeated the experiment, but this time with one agent initialized at half charge (50\%) while all others remained at 100\%. As shown in Fig.~\ref{fig:boxplot} and Table \ref{tab:results_batt_aware}, the mean $U_{i,j}^\%$ for a normal agent was \rev{60.0\%}, whereas an agent initialized with reduced battery travelled through regions with a lower mean $U_{i,j}^\%$ of \rev{49.0\%}. A t-test confirmed the statistical significance of this difference (\rev{$t = 6.54$, $p=5.0\times10^{-10}$}, $df=198$). \rev{Because a simulation ends when any battery is empty, runs with the half-charged agent are shorter (141\,s vs.\ 169\,s on average), so the two averages cover episodes of different length.} With this, we conclude that battery monitoring helps agents find adaptive strategies to increase the swarm's endurance.

\subsection{\texorpdfstring{\rev{Leadership and Position Changes}}{Leadership and Position Changes}}
\label{sec:leadership_results}
\begingroup\color{red}
We compute the metrics of Section~\ref{sec:metrics} over 30 simulations of 20 robots for each
controller (Table~\ref{tab:leadership}, Fig.~\ref{fig:leadership}). The stage-2 controller
changes its front robot about 10 times per minute, almost five times as often as the baseline
(2.1 per minute); with the safety mechanisms it does so 4.8 times per minute. Each of these
changes lasts at least 1.5\,s, so different robots genuinely take the exposed front position
in turn (Fig.~\ref{fig:leadership}a). The raw count of front changes gives the opposite
impression: the baseline's front robot changes almost 600 times per simulation, but only
about 2\% of those changes persist, because its robots move nearly level with each other and
the front identity flickers between them. For the stage-2 controller, 65\% of the raw changes
persist. For the same reason, front occupancy does not separate the controllers at 20 robots
(0.84 for stage 2 against 0.89 for the baseline): the flickering spreads the baseline's front
time over many robots. With 10 robots, where the baseline flickers less, the difference is
clear (0.90 against 0.26).

The rotation is not balanced in the sense of Voelkl et al.~\cite{voelkl2015}. The
reciprocity between front and back time is negative for both evolved controllers ($-0.26$
and $-0.19$), whereas it is close to zero for the baseline ($0.07$). Robots that spend more
time at the front spend less time at the back, so over a whole simulation some robots stay
towards the front and others towards the back, even though the front itself changes hands
frequently. None of the controllers shows a strong turning hierarchy: the leadership entropy
is at least 0.97 for all three, and the evolved controllers' hierarchies are only slightly
steeper than the baseline's (0.63 and 0.62 against 0.55). The evolved behaviour is therefore
best described as frequent exchange of the front position without fixed leaders, rather than
as a strict rotation through all positions.
\endgroup

\begin{table}[t]
\centering
\caption{\rev{Leadership metrics for 20 robots, median [interquartile range] over 30 simulations. All differences from the baseline are significant (Mann--Whitney U, $p < 0.004$).}}
\label{tab:leadership}
{\color{red}
\begin{tabular}{lccc}
\toprule
\textbf{Metric} & \textbf{Stage 2} & \textbf{Stage 2, clamped} & \textbf{Baseline} \\
\midrule
Front changes per minute & 9.9 [9.5, 10.2] & 4.8 [3.5, 5.6] & 2.1 [1.7, 2.5] \\
Raw changes that persist [\%] & 65 [57, 68] & 38 [32, 45] & 2.1 [1.8, 2.6] \\
Front occupancy entropy & 0.84 [0.81, 0.88] & 0.42 [0.36, 0.51] & 0.89 [0.88, 0.91] \\
Reciprocity $r$ & $-0.26$ [$-0.34$, $-0.21$] & $-0.19$ [$-0.23$, $-0.14$] & 0.07 [$-0.04$, 0.19] \\
Hierarchy entropy & 0.98 [0.97, 0.98] & 0.98 [0.98, 0.99] & 0.99 [0.99, 0.99] \\
Hierarchy steepness & 0.63 [0.55, 0.69] & 0.62 [0.54, 0.67] & 0.55 [0.50, 0.59] \\
\bottomrule
\end{tabular}}
\end{table}

\begin{figure}[t]
  \centering
  \includegraphics[width=\linewidth]{Images/results/paper_leadership_n20}
  \caption{\rev{Front-position changes with 20 robots. (a, b) Which robot is at the front over time in one simulation, after discarding spells shorter than 1.5\,s, for the stage-2 controller and the baseline; note the different time axes. (c) Front-position changes per minute and (d) share of raw front changes that persist for at least 1.5\,s, over 30 simulations per controller.}}
  \label{fig:leadership}
\end{figure}

\subsection{\texorpdfstring{\rev{Real-Robot Experiments}}{Real-Robot Experiments}}
\label{sec:hardware_results}
\rev{In simulation, the controller does not outperform the baseline in swarms as small as those we can run on real robots (Fig.~\ref{fig:swarm_size}). The real-robot experiments therefore test whether the evolved controller transfers from simulation to hardware, not whether it beats the baseline.}
\todo[inline]{Real-robot results: trajectories, virtual-battery traces, position changes, comparison with simulation at the same swarm size.}


\begin{figure}[t]
  \centering
  \includegraphics[width=\linewidth]{Images/results/n_agents_sweep_distance_battery_plain_seed123}
  \caption{\rev{Distance (left) and remaining battery (right) against swarm size, mean $\pm$ one standard deviation over 30 simulations per size, for the baseline (green) and the stage-2 controller evolved without (orange) and with (blue) the safety mechanisms. Both controllers were evolved with 20 robots.}}
  \label{fig:swarm_size}
\end{figure}


\begin{figure}[t]
  \centering
  \includegraphics[width=\linewidth]{Images/results/paper_trajectories_stage1_vs_stage2}
  \caption{\rev{Trajectories of the best controller from stage 1 (left) and stage 2 (right) with one (top) and five (bottom) agents; black dots mark the starting positions. Note the different axis scales in the top-right panel.}}
  \label{fig:Debunk_Compass}
\end{figure}

The source code and the supplementary video are available in the Zenodo repository at \url{https://doi.org/10.5281/zenodo.18602800}.



\section{Conclusion}
\label{sec:Conclusion}

We studied energy-efficient collective motion in a swarm of simple, memory-less ground robots. The swarm demonstrated emergent formation and dynamic formation reconfiguration, without these behaviours being explicitly encoded in the fitness function. Using a neural network controller with standard Hebbian plasticity, whose parameters were evolved via CMA-ES, the swarm learned to exploit local sensory readings and internal battery monitoring without explicit communication or wind sensing. \rev{With 20 robots, the resulting controller travelled $127\%$ further on average than a flocking baseline while also keeping more battery, and a version evolved with a speed clamp for safe use on real robots still travelled $59\%$ further. This advantage only appears in large swarms: with 12 or fewer robots the baseline performs better, and a single robot with the evolved controller hardly moves.} \rev{In the evolved swarm the front position changes hands about ten times per minute, five times as often as in the baseline, although robots do not cycle evenly between front and back.} A dedicated battery-awareness test showed that low-charge robots adapt by taking a conservative strategy through shielding themselves with other robots. Future works will investigate the minimal sensing and controller requirements that can manifest similar results\rev{, and how to evolve controllers that work across swarm sizes, for example by evaluating each candidate on several swarm sizes during evolution}. In addition, the impact of battery awareness on the learned strategies needs to be further investigated. For this matter, it is necessary to evolve a swarm that does not incorporate battery monitoring and see if similar behaviours can emerge.
